\section{The AGN-hosted fraction is measurable across its whole range}
\label{sec:fraction}

On complete clustered catalogs, \GlassNobs\ dark sirens recover the AGN-hosted
fraction everywhere between almost none and all of the population. Scanning
\fagn\ at the true expansion rate over \BaseFGridN\ points gives medians of
\BaseFMedNought, \BaseFMedThree, \BaseFMedSeven\ and \BaseFMedUnity\ against
planted values of \BaseFTruthNought, \BaseFTruthThree, \BaseFTruthSeven\ and
\BaseFTruthUnity\ (Table~\ref{tab:frecovery}, Figure~\ref{fig:frecovery}). The
planted value lies inside the 68\% interval at every interior point; the
$\fagn = 1$ case sits on the prior boundary, where an equal-tailed interval
cannot cover it and the one-sided 68\% limit, \BaseFOneSidedUnity, does. The
residual offsets are \BaseFOffNought, \BaseFOffThree, \BaseFOffSeven\ and
\BaseFOffUnity, none of them larger than \BaseFOffMaxAbs\ and all comparable to
the 68\% half-widths, which run from \BaseFHwMin\ to \BaseFHwMax.

What makes this work is the contrast between the two host priors and not the
sparse tracer's redshifts on their own. The AGN catalog contains \GlassNagn\
objects against \GlassNgal\ galaxies and leaves \GlassAgnEmptyPix\% of the sky
without a single catalogued AGN, so on its own it localises almost nothing; what
it does do is make a large part of the sky implausible for an AGN-hosted event
while the galaxy catalog covers all of it. The mixture weight is the parameter
that measures how often events prefer the sparse tracer's structure, and at
these densities it is measured to \BaseFHwMin--\BaseFHwMax\ from \GlassNobs\
events.

\begin{figure*}
\centering
\includegraphics[width=\textwidth]{figures/fig_f_recovery.pdf}
\caption{Recovery of the AGN-hosted fraction from \GlassNobs\ dark sirens on
complete clustered catalogs. \emph{(a)} the flat-prior posteriors from
\BaseFGridN-point scans at the true expansion rate, one per planted fraction
(triangles). \emph{(b)} recovered against planted, with 68\% intervals; the
$\fagn = 1$ point lies on the prior boundary and carries its one-sided limit.
No likelihood evaluation is inadmissible in any of these scans.
\label{fig:frecovery}}
\end{figure*}

\begin{deluxetable}{cccc}
\tablecaption{The AGN-hosted fraction recovered from the clustered two-tracer mock at the true expansion rate, on complete catalogs. \label{tab:frecovery}}
\tablehead{\colhead{planted $f_{\rm AGN}$} & \colhead{median} & \colhead{68\% interval} & \colhead{offset}}
\startdata
0.0099 & 0.0195 & $[0.006,\,0.039]$ & $+0.010$ \\
0.3070 & 0.3221 & $[0.300,\,0.346]$ & $+0.015$ \\
0.7030 & 0.6872 & $[0.662,\,0.712]$ & $-0.016$ \\
1.0000 & 0.9750 & $[0.966,\,1]^{\dagger}$ & $-0.025$ \\
\enddata
\tablecomments{$^{\dagger}$ At $f_{\rm AGN}=1$ the truth lies on the prior boundary, where an equal-tailed interval cannot cover it; the one-sided 68\% lower limit is quoted instead. Each row is a \BaseFGridN-point scan of \GlassNobs\ events.}
\end{deluxetable}


\subsection{The joint plane, and a distance-scale deficit}
\label{sec:fraction:joint}

Freeing the expansion rate together with the fraction leaves the fraction
essentially unchanged and exposes a systematic in the distance scale
(Figure~\ref{fig:joint}). At a planted $\fagn = \BaseFTruthThree$ the joint grid
gives $\Hz = \BaseJointHzeroThree^{+\BaseJointHzeroUpThree}_{-\BaseJointHzeroDnThree}$~\kmsmpc\
and $\fagn = \BaseJointFThree \pm \BaseJointFHwThree$ with a correlation
coefficient of $\rho = \BaseJointRhoThree$; at $\fagn = \BaseFTruthSeven$ it
gives $\Hz = \BaseJointHzeroSeven^{+\BaseJointHzeroUpSeven}_{-\BaseJointHzeroDnSeven}$~\kmsmpc,
$\fagn = \BaseJointFSeven \pm \BaseJointFHwSeven$ and
$\rho = \BaseJointRhoSeven$. The near-zero correlation at the lower fraction is
worth stating on its own: there, the mixture weight is measured essentially
independently of the distance scale, so the two parameters are not competing for
the same information.

The expansion rate, however, is recovered low --- by
\BaseJointHzeroOffAbsThree~\kmsmpc\ at $\fagn = \BaseFTruthThree$ and
\BaseJointHzeroOffAbsSeven~\kmsmpc\ at $\fagn = \BaseFTruthSeven$, which is
\BaseJointHzeroSigThree\ and \BaseJointHzeroSigSeven\ times the respective 68\%
half-widths --- and the deficit grows as the mixture leans on the sparse tracer.
It is not a property of the mixture weight: \fagn\ is recovered correctly both
at fixed \Hz\ and along the joint surface, and $\rho \simeq 0$ at the lower
fraction. Two explanations that suggest themselves are measurably wrong. The
host redshift weighting is one: the events' hosts were drawn uniformly over
eligible catalog entries, corresponding to a rate index one unit away from the
value the likelihood assumes, and repeating the scans at the matching value
moves \Hz\ by at most \BaseGammaShiftMax~\kmsmpc, because under the global sky
normalisation the factor enters both the per-event host prior and the
survey-wide normalisation and cancels. The catalog's redshift boundary is the
other, and it is not innocent: truncating both catalogs at $z \le 1$ with the
events held fixed moves \Hz\ by \BaseEdgeShiftThree~\kmsmpc\ at
$\fagn = \BaseFTruthThree$. A \BaseEdgeShiftAbsThree~\kmsmpc\ shift from
relocating a catalog-construction boundary, with the data unchanged, says the
inference is sensitive to how the host prior is terminated well outside the
range the events occupy; \S\ref{sec:selection:tilt} explains why.

Splitting the likelihood at fixed mixture weight into its per-event and
selection terms localises where that sensitivity lives. At
$\fagn = \BaseFTruthThree$ the per-event term alone peaks at
\BaseDecompNumThree~\kmsmpc\ and the selection term shifts the total by
\BaseDecompShiftThree~\kmsmpc, landing at \BaseDecompTotalThree; at
$\fagn = \BaseFTruthSeven$ the same numbers are \BaseDecompNumSeven,
\BaseDecompShiftSeven\ and \BaseDecompTotalSeven. The selection term supplies
$\sim \BaseDecompShiftAbsThree$~\kmsmpc\ of pull almost independently of the
fraction, with $\ln\mu$ spanning \BaseDecompLnmuRangeThree\ nats over the
scanned range against $\Nobs = \GlassNobs$, a lever of
$\sim \BaseDecompLeverThree$~nats; the per-event term carries all of the
$f$-dependence. This localises the sensitivity but does not by itself identify
an error, because the two terms are meant to oppose one another and only their
sum has to be unbiased. The accurate-looking result at
$\fagn = \BaseFTruthThree$ is accurate by cancellation, not by correctness:
correcting only the selection normalisation would move it from
\BaseJointHzeroOffThree\ to \BaseDecompNumOffThree~\kmsmpc.

Section~\ref{sec:selection:tilt} resolves both terms: the two pulls are
spurious levers sprung by a single inconsistency between the mocks'
true-redshift detection cut and the likelihood's distance-space selection and
host prior, and an estimator matched to the detection rule recovers the truth
at both planted fractions. The $f$-dependence visible above is the collapse of
the numerator's pull as the mixture weight moves onto the sparse tracer
(\S\ref{sec:selection:anchor}). This family of offsets is specific to the
clustered mocks' detection rule; the matched mocks of
\S\ref{sec:selection}--\ref{sec:scatter} carry no analogue of it, their own
error budget is measured and closed in \S\ref{sec:budget}, and the absolute
recoveries they quote sit on the measured estimator-overhead floor
(\OrcPairedOff$\pm$\OrcPairedSem~\kmsmpc) rather than on an unexplained
offset.

\begin{figure*}
\centering
\includegraphics[width=\textwidth]{figures/fig_joint_h0f.pdf}
\caption{Joint constraint on the expansion rate and the AGN-hosted fraction
from \GlassNobs\ dark sirens on complete clustered catalogs, at two planted
fractions. \emph{(a)} 68\% and 90\% credible regions on shared axes; crosses
mark the truth. \emph{(b)} the marginal distance-scale posteriors. The fraction
is recovered at both values while the expansion rate is low by
\BaseJointHzeroOffAbsThree\ and \BaseJointHzeroOffAbsSeven~\kmsmpc\
respectively, the deficit growing as the mixture leans on the sparse tracer.
\label{fig:joint}}
\end{figure*}
